% The Black-Litterman Model — research-note PDF.
%
% PROSE, CODE AND TABLES ARE VERBATIM from the website article:
%   site/src/content/articles/black-litterman.tsx
%   site/src/content/articles/bl-code.ts
% Metadata from lib/articles.ts; project-card fallbacks from lib/topics.ts.
%
% CHARTS generated by figures/make_black_litterman.py from the article's own
% data module — the numbers the site's <Bar> components render. That script
% records a mismatch between the module and the article's tables; both are
% reproduced as they appear rather than reconciled.
\documentclass{pyportfolios-note}

\noteshorttitle{The Black-Litterman Model}

\begin{document}
\thispagestyle{notecover}

\notemasthead{Quantitative Insights}{Tutorial}{Q3 2026}{}

\notetitle{The Black-Litterman\\Model}

\notedeck{Start from the portfolio the market already holds, then tilt it only
where you have a genuine view.}

\notelede{Plain mean-variance optimization has a dirty secret: feed it your
historical-average returns and it hands back wild, concentrated, unstable
portfolios (see the previous tutorial). Fischer Black and Robert Litterman's
1990 fix at Goldman Sachs was elegant — \textbf{start from the portfolio the
market is already holding}, then tilt it only where you have a genuine view. The
result is stable, intuitive weights that collapse back to the market when you
stay quiet.}

\notecard
  {Portfolio Optimization}
  {PyPortfolioOpt · NumPy · Pandas}
  {EWJ, EWG, EWU, EWA, EWC (iShares country ETFs)}
  {Jan 2015 – Dec 2024}
  {Turning macro views into a coherent portfolio without the weights exploding}

% ============================================================= 1 Summary =====
\noteneedroom{150bp}
\notesection{1.\notenumsep Summary}

\begin{multicols}{2}
Black-Litterman treats the \textbf{market-cap portfolio} as a rational starting
point and works \textit{backwards} to infer the expected returns that would
justify it — the ``implied equilibrium returns''. You then express views
(``German equities will outperform Japanese by 3\%'') with a confidence level,
and the model \textbf{Bayesian-blends} your views with the equilibrium prior.
Feed the blended returns into a standard optimizer and you get sensible weights
that tilt toward your views in proportion to your confidence — and revert to the
market when you have none.
\end{multicols}

% =========================================================== 2 Intuition =====
\noteneedroom{170bp}
\notesection{2.\notenumsep Intuition}

\begin{multicols}{2}
\notesubsection{2.1\notenumsep The market already did the optimization}
If everyone holds the market-cap portfolio, then — under CAPM — that portfolio is
already mean-variance efficient for \textit{some} set of expected returns.
\textbf{Reverse optimization} recovers those returns from the observed market
weights and covariance. This is the prior: not a guess, but the crowd's
collective bet.

\notesubsection{2.2\notenumsep Views as evidence, not commandments}
Rather than overwriting returns, you state a view and a \textbf{confidence}. A
low-confidence view nudges the portfolio slightly; a high-confidence view moves
it a lot. The math is Bayesian updating: prior (market) + likelihood (your view)
$\rightarrow$ posterior (blended returns).

\notesubsection{2.3\notenumsep Why the weights behave}
Because you start from the market and tilt, the output never explodes into the
200\%-long/150\%-short monsters raw MVO produces. No view on an asset? Its
weight stays at the market weight. This single property is why Black-Litterman
is the allocation workhorse of real money managers.
\end{multicols}

% ================================================= 3 Theory & Mechanics ======
\noteneedroom{190bp}
\notesection{3.\notenumsep Theory \& Mechanics}

\textbf{Step 1 — implied equilibrium returns} from market weights $w_{mkt}$ and
covariance $\Sigma$:

\[ \Pi = \delta\, \Sigma\, w_{mkt} \]

where $\delta$ is the market risk-aversion coefficient.

\textbf{Step 2 — the views}, encoded as $P\,\mu = Q + \varepsilon$, with $P$
picking the assets, $Q$ the view magnitudes, and $\Omega$ the view uncertainty.

\textbf{Step 3 — the Black-Litterman posterior returns:}

\[ \mu_{BL} = \Big[(\tau\Sigma)^{-1} + P^\top\Omega^{-1}P\Big]^{-1}
   \Big[(\tau\Sigma)^{-1}\Pi + P^\top\Omega^{-1}Q\Big] \]

$\tau$ scales the uncertainty of the prior. We'll let \textbf{PyPortfolioOpt}
handle the linear algebra and focus on what goes in and what comes out.

% ============================== 4 Applied Example — Five Country ETFs ========
\noteneedroom{200bp}
\notesection{4.\notenumsep Applied Example — Five Country ETFs}

\notesubsection{4.1\notenumsep A global equity universe}

Five developed markets via iShares MSCI country ETFs — the natural setting for
Black-Litterman, where ``views'' are macro calls on countries:

\begin{notetable}{@{}p{110bp}p{386bp}@{}}
\thcell{Ticker} & \thcell{Country} \\
\theadrule
EWJ & Japan \\ \rowrule
EWG & Germany \\ \rowrule
EWU & United Kingdom \\ \rowrule
EWA & Australia \\ \rowrule
EWC & Canada \\
\end{notetable}

\begin{notecode}
TICKERS = ["EWJ", "EWG", "EWU", "EWA", "EWC"]
COUNTRY = {"EWJ": "Japan", "EWG": "Germany", "EWU": "UK", "EWA": "Australia", "EWC": "Canada"}
START, END = "2015-01-01", "2024-12-31"

def load_prices(tickers, start, end):
    """Adjusted-close prices: yfinance first, Stooq as fallback."""
    try:
        import yfinance as yf
        df = yf.download(tickers, start=start, end=end, auto_adjust=True, progress=False)["Close"]
        if not df.empty:
            return df[tickers].dropna()
    except Exception as exc:
        print(f"yfinance failed ({exc}); trying Stooq…")
    cols = {}
    for t in tickers:
        url = f"https://stooq.com/q/d/l/?s={t.lower()}.us&i=d"
        cols[t] = pd.read_csv(url, parse_dates=["Date"], index_col="Date")["Close"].rename(t)
    return pd.concat(cols, axis=1).loc[start:end].dropna()

px = load_prices(TICKERS, START, END)
rets = px.pct_change().dropna()
print(f"{len(px)} trading days, {px.index[0].date()} → {px.index[-1].date()}")
\end{notecode}

\notesubsection{4.2\notenumsep Market weights and implied equilibrium returns}

Black-Litterman needs market-cap weights as the prior. Ideally these come from
each market's investable cap; here we use approximate relative market sizes as a
stand-in (in production you'd pull real market caps). From these weights we
reverse-engineer the returns that justify them.

\begin{notecode}
from pypfopt import risk_models, expected_returns
from pypfopt import black_litterman
from pypfopt.black_litterman import BlackLittermanModel

S = risk_models.sample_cov(px)

# approximate relative market caps (illustrative stand-in for true investable caps)
mcaps = {"EWJ": 6.0, "EWG": 2.5, "EWU": 3.0, "EWA": 1.6, "EWC": 2.8}   # in trillions USD, rough
w_mkt = pd.Series(mcaps) / sum(mcaps.values())

delta = 2.5   # market risk-aversion (typical value)
pi = black_litterman.market_implied_prior_returns(mcaps, delta, S)

prior_tbl = pd.DataFrame({"market weight": w_mkt, "implied return": pi}).round(3)
prior_tbl.index = [f"{t} ({COUNTRY[t]})" for t in prior_tbl.index]
print(prior_tbl)
\end{notecode}

\begin{notetable}{@{}p{200bp}>{\raggedleft\arraybackslash}p{140bp}>{\raggedleft\arraybackslash}p{140bp}@{}}
\thcell{Market} & \thcell{Market weight} & \thcell{Implied return} \\
\theadrule
EWJ (Japan) & 37.7\% & 6.6\% \\ \rowrule
EWG (Germany) & 15.7\% & 8.2\% \\ \rowrule
EWU (UK) & 18.9\% & 7.6\% \\ \rowrule
EWA (Australia) & 10.1\% & 8.8\% \\ \rowrule
EWC (Canada) & 17.6\% & 7.3\% \\
\end{notetable}

Read the implied returns as ``what the crowd must believe'': higher-beta markets
(Australia) need higher expected returns to justify their market weight.

\notesubsection{4.3\notenumsep Express a view}

Our macro call: \textbf{Germany (EWG) will return 10\% annually} — versus an
implied 8.2\% — and we're \textbf{moderately confident} (50\%, via Idzorek's
method). The model figures out the full portfolio implications on its own.

\begin{notecode}
# absolute + relative views
viewdict = {
    "EWG": 0.10,                 # absolute: Germany returns 10%
}
# a relative view via P/Q would go here too; keep one absolute view for clarity

bl = BlackLittermanModel(S, pi=pi, absolute_views=viewdict, omega="idzorek",
                         view_confidences=[0.50])
bl_returns = bl.bl_returns()

compare = pd.DataFrame({"implied (prior)": pi, "Black-Litterman (posterior)": bl_returns}).round(3)
compare.index = [f"{t} ({COUNTRY[t]})" for t in compare.index]
print(compare)
print("\nNote: only EWG had a view — but the posterior nudges correlated markets too.")
\end{notecode}

\begin{notetable}{@{}p{200bp}>{\raggedleft\arraybackslash}p{140bp}>{\raggedleft\arraybackslash}p{140bp}@{}}
\thcell{Market} & \thcell{Implied (prior)} & \thcell{Posterior (with view)} \\
\theadrule
EWJ (Japan) & 6.6\% & 7.1\% \\ \rowrule
EWG (Germany) & 8.2\% & 9.1\% \\ \rowrule
EWU (UK) & 7.6\% & 8.3\% \\ \rowrule
EWA (Australia) & 8.8\% & 9.6\% \\ \rowrule
EWC (Canada) & 7.3\% & 8.0\% \\
\end{notetable}

\notefigure{figures/black-litterman/fig1-returns.png}
  {Figure 4.3 · One view on Germany moves every posterior — via correlation}
  {Implied prior against Black-Litterman posterior; Germany carries the view,
   but every correlated market shifts too}
  {Source: Author calculations, from the article's computed results.}

Only EWG had a view — but every posterior moved, because the markets are
correlated. A bullish call on Germany is implicitly (weaker) good news for
everything that co-moves with Germany.

\notesubsection{4.4\notenumsep Optimize on the blended returns}

Feed the posterior returns into a max-Sharpe optimizer and compare the
Black-Litterman weights against both the market prior and what naive MVO on
historical means would have produced.

\begin{notecode}
from pypfopt import EfficientFrontier

# Black-Litterman optimal weights
ef_bl = EfficientFrontier(bl_returns, S)
ef_bl.max_sharpe(risk_free_rate=0.02)
w_bl = pd.Series(ef_bl.clean_weights())

# naive MVO on historical means, for contrast
mu_hist = expected_returns.mean_historical_return(px)
ef_naive = EfficientFrontier(mu_hist, S)
ef_naive.max_sharpe(risk_free_rate=0.02)
w_naive = pd.Series(ef_naive.clean_weights())

weights = pd.DataFrame({"Market prior": w_mkt, "Black-Litterman": w_bl, "Naive MVO": w_naive}).round(3)
weights.index = [COUNTRY[t] for t in weights.index]
print(weights)

ax = weights.plot(kind="bar", figsize=(10, 5), width=0.8,
                  color=["gray", "steelblue", "darkorange"])
ax.set_ylabel("Weight"); ax.set_title("Allocations: market prior vs Black-Litterman vs naive MVO")
ax.axhline(0, color="black", lw=0.6)
plt.tight_layout(); plt.show()
\end{notecode}

\begin{notetable}{@{}p{136bp}>{\raggedleft\arraybackslash}p{120bp}>{\raggedleft\arraybackslash}p{130bp}>{\raggedleft\arraybackslash}p{90bp}@{}}
\thcell{Market} & \thcell{Market prior} & \thcell{Black-Litterman} & \thcell{Naive MVO} \\
\theadrule
Japan & 37.7\% & 20.5\% & 74.8\% \\ \rowrule
Germany & 15.7\% & 35.6\% & 0.0\% \\ \rowrule
UK & 18.9\% & 10.8\% & 0.0\% \\ \rowrule
Australia & 10.1\% & 25.7\% & 0.0\% \\ \rowrule
Canada & 17.6\% & 7.4\% & 25.2\% \\
\end{notetable}

\notefigure{figures/black-litterman/fig2-weights.png}
  {Figure 4.4 · Allocations — market prior vs Black-Litterman vs naive MVO}
  {Naive MVO concentrates and takes short positions; Black-Litterman stays near
   market weights and tilts only where a view was expressed}
  {Source: Author calculations, from the article's computed results. The naive-MVO
   series here is the unconstrained solution, which is why it shows shorts.}

The story in one chart: \textbf{naive MVO} lurches to extremes — 74.8\% in Japan
(the decade's backtest winner) and \textit{zero} in three of five markets —
while \textbf{Black-Litterman} holds every market and tilts sensibly toward
Germany (15.7\% $\rightarrow$ 35.6\%) where we expressed our view. That stability
is the whole point.

\notesubsection{4.5\notenumsep Sanity check: no views $\rightarrow$ back to the market}

The definitive test. With \textit{no} views, the Black-Litterman posterior must
equal the implied prior, and the optimized weights must return to market
weights. Our run gives a maximum posterior-minus-prior difference of exactly
\textbf{0.0} — the model is genuinely neutral by default.

\begin{notecode}
bl_noview = BlackLittermanModel(S, pi=pi, absolute_views={}, omega="idzorek", view_confidences=[])
diff = (bl_noview.bl_returns() - pi).abs().max()
print(f"Max |posterior - prior| with no views: {diff:.2e}  (should be ~0)")
\end{notecode}

% ========================================================== 5 Conclusion =====
\noteneedroom{330bp}
\notesection{5.\notenumsep Conclusion}

\begin{multicols}{2}
\notesubsection{5a.\notenumsep Strengths}
\begin{notebullets}
\item \textbf{Stable, intuitive weights} — no more 200\%-long/150\%-short
      monsters from raw MVO
\item \textbf{Neutral by default} — with no views, you hold the market; tilts
      scale with confidence
\item \textbf{Combines art and science} — a formal channel for a manager's macro
      views inside a rigorous framework
\item \textbf{Partial views work} — a view on one asset sensibly adjusts the
      whole portfolio via correlations
\end{notebullets}

\notesubsection{5b.\notenumsep Weaknesses \& Limitations}
\begin{notebullets}
\item \textbf{Needs market-cap weights} — the prior is only as good as your
      market-weight and covariance inputs
\item \textbf{τ and Ω are fiddly} — the uncertainty parameters have no universal
      setting and affect the outcome
\item \textbf{Still mean-variance underneath} — inherits variance's blindness to
      tail risk
\item \textbf{Views are yours to get right} — the model propagates a wrong view
      as faithfully as a right one
\end{notebullets}

\notesubsection{5c.\notenumsep Applications in Practice}
\begin{notebullets}
\item The allocation engine at real asset managers, precisely because the
      weights are usable as-is
\item Turning a research team's country/sector calls into a coherent portfolio
\item Tactical tilts around a strategic (market) benchmark
\end{notebullets}

\notesubsection{5d.\notenumsep Alternatives \& Extensions}
\begin{notebullets}
\item \textbf{Idzorek's confidence method} — specify view confidence as an
      intuitive 0–100\% instead of an abstract Ω (used above)
\item \textbf{Risk parity} — skips return estimation entirely (the next tutorial)
\item \textbf{Entropy pooling (Meucci)} — a more general way to impose views,
      including on higher moments
\item \textbf{Mean-variance optimization} — the raw method Black-Litterman was
      invented to tame (the previous tutorial)
\end{notebullets}
\end{multicols}

\end{document}
